% concatenated from lebesgue_constant_v4_sn.tex
%Version 3.1 December 2024
% See section 11 of the User Manual for version history
%
%%%%%%%%%%%%%%%%%%%%%%%%%%%%%%%%%%%%%%%%%%%%%%%%%%%%%%%%%%%%%%%%%%%%%%
%%                                                                 %%
%% Please do not use \input{...} to include other tex files.       %%
%% Submit your LaTeX manuscript as one .tex document.              %%
%%                                                                 %%
%% All additional figures and files should be attached             %%
%% separately and not embedded in the \TeX\ document itself.       %%
%%                                                                 %%
%%%%%%%%%%%%%%%%%%%%%%%%%%%%%%%%%%%%%%%%%%%%%%%%%%%%%%%%%%%%%%%%%%%%%

%%\documentclass[referee,sn-basic]{sn-jnl}% referee option is meant for double line spacing

%%=======================================================%%
%% to print line numbers in the margin use lineno option %%
%%=======================================================%%

%%\documentclass[lineno,pdflatex,sn-basic]{sn-jnl}% Basic Springer Nature Reference Style/Chemistry Reference Style

%%=========================================================================================%%
%% the documentclass is set to pdflatex as default. You can delete it if not appropriate.  %%
%%=========================================================================================%%

%%\documentclass[sn-basic]{sn-jnl}% Basic Springer Nature Reference Style/Chemistry Reference Style

%%Note: the following reference styles support Namedate and Numbered referencing. By default the style follows the most common style. To switch between the options you can add or remove ---Numbered---?in the optional parenthesis. 
%%The option is available for: sn-basic.bst, sn-chicago.bst%  
 
%%\documentclass[pdflatex,sn-nature]{sn-jnl}% Style for submissions to Nature Portfolio journals
%%\documentclass[pdflatex,sn-basic]{sn-jnl}% Basic Springer Nature Reference Style/Chemistry Reference Style
\documentclass[pdflatex,sn-mathphys-num]{sn-jnl}% Math and Physical Sciences Numbered Reference Style
%%\documentclass[pdflatex,sn-mathphys-ay]{sn-jnl}% Math and Physical Sciences Author Year Reference Style
%%\documentclass[pdflatex,sn-aps]{sn-jnl}% American Physical Society (APS) Reference Style
%%\documentclass[pdflatex,sn-vancouver-num]{sn-jnl}% Vancouver Numbered Reference Style
%%\documentclass[pdflatex,sn-vancouver-ay]{sn-jnl}% Vancouver Author Year Reference Style
%%\documentclass[pdflatex,sn-apa]{sn-jnl}% APA Reference Style
%%\documentclass[pdflatex,sn-chicago]{sn-jnl}% Chicago-based Humanities Reference Style

%%%% Standard Packages
%%<additional latex packages if required can be included here>

\usepackage{graphicx}%
\usepackage{multirow}%
\usepackage{amsmath,amssymb,amsfonts}%
\usepackage{amsthm}%
\usepackage{mathrsfs}%
\usepackage[title]{appendix}%
\usepackage{xcolor}%
\usepackage{textcomp}%
\usepackage{manyfoot}%
\usepackage{booktabs}%
\usepackage{algorithm} % ADD THIS PACKAGE! It provides the 'algorithm' environment.
\usepackage{algorithmicx}% This is needed by algpseudocode.
\usepackage{algpseudocode}% This provides the pseudocode commands.
\usepackage{listings}%
\usepackage{bm}
\usepackage{subcaption}
%%%%

%%%%%=============================================================================%%%%
%%%%  Remarks: This template is provided to aid authors with the preparation
%%%%  of original research articles intended for submission to journals published 
%%%%  by Springer Nature. The guidance has been prepared in partnership with 
%%%%  production teams to conform to Springer Nature technical requirements. 
%%%%  Editorial and presentation requirements differ among journal portfolios and 
%%%%  research disciplines. You may find sections in this template are irrelevant 
%%%%  to your work and are empowered to omit any such section if allowed by the 
%%%%  journal you intend to submit to. The submission guidelines and policies 
%%%%  of the journal take precedence. A detailed User Manual is available in the 
%%%%  template package for technical guidance.
%%%%%=============================================================================%%%%

%% as per the requirement new theorem styles can be included as shown below
\theoremstyle{thmstyleone}%
\newtheorem{definition}{Definition}
\newtheorem{theorem}{Theorem}
\newtheorem{lemma}{Lemma}
\newtheorem{remark}{Remark}
\newtheorem{corollary}{Corollary}

\newcommand{\norm}[1]{\left\lVert#1\right\rVert}
\newcommand{\abs}[1]{\left\lvert#1\right\rvert}
\newcommand{\N}{\mathbb{N}}
\newcommand{\mat}[1]{\mathbf{#1}}
\newcommand{\vect}[1]{\mathbf{#1}}


\raggedbottom
%%\unnumbered% uncomment this for unnumbered level heads

\begin{document}

\title[On the Pre-Asymptotic Stability and Inverse Structure of Extended-Domain Spectral Methods]{On the Pre-Asymptotic Stability and Inverse Structure of Extended-Domain Spectral Methods}

%%=============================================================%%
%% GivenName	-> \fnm{Joergen W.}
%% Particle	-> \spfx{van der} -> surname prefix
%% FamilyName	-> \sur{Ploeg}
%% Suffix	-> \sfx{IV}
%% \author*[1,2]{\fnm{Joergen W.} \spfx{van der} \sur{Ploeg} 
%%  \sfx{IV}}\email{iauthor@gmail.com}
%%=============================================================%%

\author*[1]{\fnm{Po-Yi} \sur{Wu}}\email{r04527030@gmail.com}



\affil*[1]{Institute of Applied Mechanics, National Taiwan University, No. 1, Sec. 4, Roosevelt Rd., 106319,  Taipei, Taiwan}

\abstract{The extended-domain method is a strategy for applying spectral methods to complex geometries. Its stability is complicated by the ill-conditioning of the Fourier extension frame. This paper provides an analysis of the method's pre-asymptotic behavior. We confirm that the spectral collocation system is asymptotically ill-conditioned for both the Poisson and convection-diffusion operators, driven by the redundancy of the underlying frame. However, we prove a fundamental structural dichotomy in their discrete Green's functions. We show that the inverse of the convection-diffusion operator is numerically highly asymmetric, exhibiting exponential upstream decay, in contrast to the numerically dense inverse of the Poisson operator. This intrinsic asymmetry explains why the convection-diffusion operator is significantly more robust to the underlying frame instability in practical computations.} 
\keywords{Spectral method, Lebesgue constant, extended domain, convection-diffusion, Poisson equation, stability analysis. MSC Classification: 65N12, 65N35, 65T40}


\maketitle
%%==================================%%
%% Sample for unstructured abstract %%
%%==================================%%


\section{Introduction}
%================================================================================

Spectral methods are a cornerstone of scientific computing, widely used for their ability to achieve exponential convergence rates for smooth functions \cite{boyd2001chebyshev, canuto2007spectral}. However, their classical formulation relies heavily on tensor-product grids, making their application to complex, non-rectangular geometries challenging. To bridge this gap, various ``fictitious domain'' strategies have been proposed, including interface tracking \cite{leveque1994immersed}, penalization methods \cite{angot1999penalization}, and smooth extension techniques \cite{bruno2010high, stein2016immersed}.

% 

Among these strategies, the extended-domain method (also known as Fourier extension) is conceptually attractive due to its simplicity. The method embeds the complex physical domain within a larger, regular computational box and approximates the solution using a Fourier series defined on this extended domain \cite{lui2009spectral, sabetghadam2009fourier, fang2011towards}.

Despite its utility, the numerical stability of the extended-domain method is a subject of nuance and ongoing research. The fundamental difficulty lies in the approximation properties of the basis. When a Fourier series on an extended domain is restricted to a subset (the physical domain), the basis functions form a frame rather than an orthonormal basis. This redundancy leads to an interpolation problem that is highly ill-conditioned. Indeed, the seminal work of Platte et al. \cite{platte2011impossibility} established that for equispaced data, no method can simultaneously achieve geometric convergence and stability; the condition number must grow exponentially if the convergence is geometric. To mitigate this, standard practice often employs oversampling (where the number of collocation points $M$ exceeds the number of modes $N$) or regularization to stabilize the frame \cite{huybrechs2010fourier, adcock2014numerical}.


However, a discrepancy persists in the pre-asymptotic regime, particularly when the system is square ($M=N$) or only mildly oversampled. Numerical evidence suggests that the method is more robust for the non-self-adjoint convection-diffusion equation than for the self-adjoint Poisson equation. While the underlying basis is unstable in both cases, the convection-diffusion operator appears to suppress this instability in practical computations, whereas the Poisson operator does not. This resonates with the observations of Trefethen et al. \cite{Trefethen1999} regarding non-normal operators, where standard eigenvalue analysis often fails to predict transient behavior.

This paper provides a structural analysis that resolves this situation. We investigate the spectral collocation method in the ``square'' limit ($M=N$) to isolate the interaction between the differential operator and the basis ill-conditioning. While regularization is necessary for practical solvers, analyzing the square limit acts as a magnifying glass, revealing the intrinsic structural differences in how the operators propagate the frame's inherent instability. Our contributions are as follows:
\begin{itemize}
    \item We characterize the method's baseline instability. Consistent with the theory of Fourier frames \cite{adcock2014numerical, platte2011impossibility}, we demonstrate that the condition number of the linear system grows exponentially. This confirms that the ultimate source of instability is the redundancy of the basis, which affects all operators asymptotically.
    
    \item We explain the observed pre-asymptotic stability of the convection-diffusion operator by proving a fundamental structural dichotomy in the discrete Green's functions (the inverses of the physical-space operators). We establish that the inverse of the discrete convection-diffusion operator is \textbf{numerically highly asymmetric}, exhibiting exponential upstream decay (visualized later in Figure~\ref{fig:matrix_structures}). This is in contrast to the \textbf{numerically dense} inverse of the Poisson operator.
    
    \item We demonstrate that this structural dichotomy dictates the stability of the unstabilized square system ($M=N$). While the discrete system lacks the coercivity required for sharp entry-wise bounds in the absence of stabilization, we establish that the continuous physical operator strictly attenuates boundary-localized frame errors. This physical attenuation competes against the exponentially growing Lebesgue constant, leading to an explicit characterization of the pre-asymptotic regime: the interior solution remains stable as long as the physical decay rate exceeds the frame growth rate. We identify the \textbf{modal P\'eclet number} as the mechanism governing the eventual transition to asymptotic instability.
\end{itemize}

These results provide a theoretical foundation that reconciles the method's known asymptotic limitations with its practical effectiveness in convection-dominated regimes.

It is important to acknowledge that Fourier extension methods, when equipped with appropriate stabilizations---such as the AZ algorithm \cite{coppe2020az}, carefully chosen regularization parameters, or optimal oversampling---have proven effective and robust in practice for a wide variety of PDEs. The purpose of this paper is not to suggest that the method is impractical. Rather, our goal is to provide a \emph{structural} explanation for why certain operators inherently resist the frame's ill-conditioning better than others \emph{before} these external stabilizations are applied. This analysis complements the practical literature by identifying which problems benefit most from the natural operator structure and which require more aggressive stabilization.

%================================================================================
\section{Spectral Collocation on Extended Domains and Baseline Instability}
%================================================================================

We consider a one-dimensional linear elliptic operator $\mathcal{L}$ on a physical domain $I = (0, L)$. The core idea of the extended-domain method is to solve the problem on $I$ using a basis of functions defined on a larger, computationally convenient domain $\tilde{I} = (-\delta, L+\delta)$ for some extension parameter $\delta > 0$.

%--------------------------------------------------------------------------------
\subsection{Basis and Discretization}
%--------------------------------------------------------------------------------

The method is built upon a specific choice of basis functions on this extended domain.

\begin{definition}[Laplacian Eigenbasis on the Extended Domain]
The basis functions $\{w_j(x)\}_{j=1}^\infty$ are the normalized eigenfunctions of the negative Laplacian on the extended domain $\tilde{I}$ with homogeneous Dirichlet boundary conditions:
\begin{equation}
-w'' = \lambda_j w \quad \text{in } \tilde{I}, \quad w(-\delta) = w(L+\delta) = 0.
\end{equation}
The eigenfunctions and eigenvalues are given by
\begin{equation} \label{eq:basis_func}
w_j(x) = \sqrt{\frac{2}{L+2\delta}} \sin\left(\frac{j\pi(x+\delta)}{L+2\delta}\right), \quad \lambda_j = \left(\frac{j\pi}{L+2\delta}\right)^2, \quad j \in \N.
\end{equation}
\end{definition}

\begin{figure}[ht!]
    \centering
    \includegraphics[width=0.9\textwidth]{extended_domain_diagram.pdf}
    \caption{Schematic of the Extended-Domain Method in 1D. The complex physical problem on $I=[0,L]$ is embedded into a larger, regular computational domain $\tilde{I}=[-\delta, L+\delta]$. The solution is approximated by Fourier-like basis functions (e.g., $w_j(x)$) that are defined on the full domain $\tilde{I}$ and vanish at the extended boundaries $-\delta$ and $L+\delta$, but generally do not satisfy boundary conditions at the physical endpoints $0$ and $L$.}
    \label{fig:domain_geometry}
\end{figure}

The geometric configuration of the physical domain $I$ embedded within $\tilde{I}$ is illustrated in Figure \ref{fig:domain_geometry}. These functions form a complete orthonormal basis in $L^{2}(\tilde{I})$. The solution to the PDE is then approximated using a truncated expansion in this basis.

\begin{definition}[Spectral Collocation Method]
The solution is approximated in the finite-dimensional space $V_N = \text{span}\{w_1, \dots, w_N\}$ by an expansion $u_N(x) = \sum_{j=1}^N c_j w_j(x)$. Given a differential operator $\mathcal{L}$, the functions $\{\phi_j(x) = \mathcal{L}[w_j](x)\}_{j=1}^N$ form the collocation basis. For a set of $N$ distinct collocation points $\{x_k\}_{k=1}^N \subset I$, the coefficients $\vect{c} = (c_1, \dots, c_N)^T$ are found by enforcing the PDE at these points:
\begin{equation}
\mathcal{L}[u_N](x_k) = \sum_{j=1}^N c_j \phi_j(x_k) = f(x_k), \quad k=1, \dots, N.
\end{equation}
This defines an $N \times N$ linear system $\mat{A}\vect{c} = \vect{f}$ where the collocation matrix entries are $(\mat{A})_{kj} = \phi_j(x_k)$.
\end{definition}

The structure of the operator-applied functions $\{\phi_j\}$ is determined entirely by the differential operator $\mathcal{L}$ and is the primary object of our stability analysis. For the two operators central to this paper, these functions are:
\begin{itemize}
    \item \textbf{Poisson Operator} ($\mathcal{L} = -d^2/dx^2$): The operator-applied functions are simply the eigenfunctions scaled by their eigenvalues, revealing a strong dependence on the mode number $j$:
    \[
    \phi_j(x) = -w_j''(x) = \lambda_j w_j(x) \propto j^2 \sin\left(\frac{j\pi(x+\delta)}{L+2\delta}\right).
    \]
    \item \textbf{Convection-Diffusion Operator} ($\mathcal{L} = -d^2/dx^2 + k(x)d/dx$): The operator-applied functions involve both second and first derivatives:
    \[
    \phi_j(x) = -w_j''(x) + k(x)w_j'(x) = \lambda_j w_j(x) + k(x)w_j'(x).
    \]
\end{itemize}
The central question of this paper is how the addition of the $k(x)w_j'(x)$ term affects the stability of the collocation scheme built upon these functions.

\begin{remark}[Independence of the Solution Basis from the Convection Coefficient]
\label{rem:basis_independence}
It is important to emphasize that the \emph{solution basis} $\{w_j\}_{j=1}^N$ is always the set of Laplacian eigenfunctions on $\tilde{I}$, defined in Eq.~\eqref{eq:basis_func}, and is entirely independent of the convection coefficient $k(x)$. The functions $\phi_j(x) = \mathcal{L}[w_j](x)$ are \emph{not} a separate basis to be chosen; they are the images of the fixed basis under the given differential operator, and they define the entries of the collocation matrix via $(\mat{A})_{kj} = \phi_j(x_k)$. In practice, the coefficient $k(x)$ is simply evaluated pointwise at the collocation nodes $\{x_k\}$; no Fourier interpolation or approximation of $k(x)$ is required. The appearance of $k(x)$ in the expression for $\phi_j$ reflects the structure of the differential operator, not a choice of basis.
\end{remark}

\begin{remark}[Enforcement of Boundary Conditions]
\label{rem:boundary_conditions}
The basis functions $\{w_j\}$ satisfy homogeneous Dirichlet conditions at the boundaries of the \emph{extended} domain $\tilde{I}$, not at the physical boundaries of $I$. To impose boundary conditions on the physical domain, we employ basis recombination: the original basis is linearly combined to produce a new basis that exactly satisfies the desired boundary values at $x=0$ and $x=L$. Algebraically, this replaces two collocation equations with boundary constraints. Importantly, this recombination operates within the same Fourier extension frame, so the asymptotic conditioning properties analyzed below are not altered.
\end{remark}
%--------------------------------------------------------------------------------
\subsection{Stability and the Lebesgue Constant}
%--------------------------------------------------------------------------------

The stability of the collocation scheme is determined by the conditioning of the matrix $\mat{A}$. A robust and basis-independent measure of this stability is the Lebesgue constant of the underlying interpolation operator.

\begin{definition}[Lebesgue Constant]
The interpolation operator $\mathcal{I}_N: C(I) \to V_N$ maps a continuous function $g$ to its unique interpolant $g_N = \mathcal{I}_N g$ in the trial space $V_N = \text{span}\{w_j\}$ that matches $g$ at the collocation points. The Lebesgue constant is the operator norm of $\mathcal{I}_N$ induced by the infinity norm:
\begin{equation}
\Lambda_N = \sup_{g \in C(I), g \ne 0} \frac{\norm{\mathcal{I}_N g}_{L^\infty(I)}}{\norm{g}_{L^\infty(I)}}.
\end{equation}
(We note that the collocation scheme also involves interpolation into the test space $\text{span}\{\phi_j\}$. Because the second-derivative term dominates for high frequencies ($\phi_j \approx \lambda_j w_j$), the Lebesgue constants for both spaces are asymptotically equivalent. We denote both generically by $\Lambda_N$, as they are driven by the exact same underlying frame redundancy.)
It can be computed via the Lagrange basis functions $\{l_k(x)\}_{k=1}^N$, which satisfy $l_k(x_j) = \delta_{kj}$, as:
\begin{equation}
\Lambda_N = \max_{x \in I} \sum_{k=1}^N \abs{l_k(x)}.
\end{equation}
\end{definition}

\begin{remark}[Lebesgue Constant as an Instability Detector]
\label{rem:lebesgue_condition_link}
The Lebesgue constant relates approximation stability to algebraic stability via a lower bound. Specifically, for the synthesis matrix $\mat{W}_N$ (defined below in Section 2.3), standard frame theory establishes that its condition number grows at least as fast as the Lebesgue constant, $\kappa_2(\mat{W}_N) \ge C \Lambda_N / \sqrt{N}$ \citep{adcock2014numerical}. As we will establish via the factorization $\mat{A}_P = \mat{W}_N\mathbf{\Lambda}$ in Section 2.3, an exponential growth in $\Lambda_N$ necessitates an exponential growth in the system condition number $\kappa(\mat{A}_P)$. Thus, while a small $\Lambda_N$ does not guarantee a well-conditioned matrix, a large $\Lambda_N$ serves as a sufficient condition to prove algebraic instability.
\end{remark}

A Lebesgue constant $\Lambda_N$ that grows polynomially with $N$ is a necessary condition for the convergence of the method, ensuring that the interpolation error does not diverge for smooth functions. In the context of this paper, we use $\Lambda_N$ to evaluate the baseline failure of the Poisson operator (Section 5.1). The algebraic stability of the Convection-Diffusion operator, which occurs despite the potential for frame ill-conditioning, is driven by the highly asymmetric inverse structure proven in Section 3.

We now establish the baseline stability properties of the extended-domain method when applied to the standard Poisson equation. We focus on the case of a square system ($M=N$), where the number of collocation points equals the number of basis modes. While oversampling is known to stabilize such frames \cite{adcock2014numerical}, the square case serves as a critical baseline to understand the inherent interaction between the basis redundancy and the differential operator.

Consistent with the impossibility results for equispaced data \cite{platte2011impossibility}, we show that the linear system for the Poisson operator is exponentially ill-conditioned. This ill-conditioning is intrinsic to the Fourier extension frame and provides the necessary context for the structural dichotomy observed in the convection-diffusion operator in the next section.

%--------------------------------------------------------------------------------
\subsection{Structure of the Collocation Matrix}
%--------------------------------------------------------------------------------

Recall that the collocation matrix $\mat{A}$ maps the coefficients $\vect{c}$ of the expansion $u_N = \sum c_j w_j$ to the values of the operator $\mathcal{L}[u_N]$ at the collocation points. For the Poisson operator $\mathcal{L} = -d^2/dx^2$, the basis functions are eigenfunctions, $\mathcal{L}[w_j] = \lambda_j w_j$. If we consider the unconstrained bulk operator (i.e., enforcing the PDE at every node without replacing rows for boundary conditions), this allows us to exactly factor the collocation matrix as:
\begin{equation}
\label{eq:poisson_factorization}
\mat{A}_P = \mat{W}_N \mathbf{\Lambda},
\end{equation}
where:
\begin{itemize}
    \item $\mathbf{\Lambda} = \text{diag}(\lambda_1, \dots, \lambda_N)$ is the diagonal matrix of eigenvalues, with $\lambda_j \propto j^2$.
    \item $\mat{W}_N$ is the $N \times N$ synthesis matrix (or frame matrix) with entries $(\mat{W}_N)_{kj} = w_j(x_k)$. It maps the coefficients to the function values on the physical grid. (We use the subscript $N$ to emphasize the dependence on the truncation level.)
\end{itemize}
While practical implementations modify this matrix to enforce boundary conditions (as noted in Remark \ref{rem:boundary_conditions}), analyzing this unconstrained bulk operator isolates the baseline exponential ill-conditioning caused by the underlying frame, as boundary constraints only further restrict the space.

The stability of the numerical solution $\vect{c} = \mat{A}_P^{-1} \vect{f}$ is governed by the condition number $\kappa(\mat{A}_P) = \norm{\mat{A}_P}_2 \norm{\mat{A}_P^{-1}}_2$.

%--------------------------------------------------------------------------------
\subsection{Exponential Growth of the Condition Number}
%--------------------------------------------------------------------------------

The instability of the method arises from the properties of the synthesis matrix $\mat{W}_N$. The set of basis functions $\{w_j\}_{j=1}^N$ restricted to the physical domain $I \subset \tilde{I}$ forms a truncated Fourier extension frame. A defining characteristic of such frames is the rapid decay of their singular values.

% 

\begin{theorem}[Exponential Ill-Conditioning of the Poisson System]
\label{thm:poisson_illcond}
Let $\mat{A}_P$ be the $N \times N$ spectral collocation matrix for the Poisson equation on the physical domain $I = (0,L) \subset \tilde{I} = (-\delta, L+\delta)$. The condition number $\kappa(\mat{A}_P)$ is bounded below by an exponential in $N$. Specifically, there exist constants $C > 0$ and $\alpha = \alpha(\delta, L) > 0$, where $\alpha$ depends on the extension ratio $\rho := L/(L+2\delta) < 1$, such that:
\begin{equation}
\kappa(\mat{A}_P) \ge C e^{\alpha(\delta, L)\, N}.
\end{equation}
In particular, $\alpha \to 0$ as $\delta \to 0$ (no extension) and $\alpha$ increases with $\delta$, reflecting the growing redundancy of the frame.
\end{theorem}

\begin{proof}
Because $\mat{A}_P$ includes the unbounded differential operator $\mathbf{\Lambda}$, its norm grows rapidly with $N$ ($\norm{\mat{A}_P}_2 \ge \lambda_N \ge 1$ for large $N$). Therefore, the condition number is strictly bounded below by the norm of the inverse: $\kappa(\mat{A}_P) = \norm{\mat{A}_P}_2 \norm{\mat{A}_P^{-1}}_2 \ge \norm{\mat{A}_P^{-1}}_2$. Using the unconstrained factorization in \eqref{eq:poisson_factorization}, the inverse is $\mat{A}_P^{-1} = \mathbf{\Lambda}^{-1} \mat{W}_N^{-1}$.
We can bound the spectral norm of the product from below:
\begin{equation}
\norm{\mat{A}_P^{-1}}_2 = \norm{\mathbf{\Lambda}^{-1} \mat{W}_N^{-1}}_2 \ge \sigma_{\min}(\mathbf{\Lambda}^{-1}) \norm{\mat{W}_N^{-1}}_2.
\end{equation}
We analyze the two terms separately:
\begin{enumerate}
    \item \textbf{The Differential Operator (Polynomial Scaling):} The matrix $\mathbf{\Lambda}^{-1}$ contains the inverse eigenvalues. Its smallest singular value corresponds to the largest eigenvalue $\lambda_N$:
    \begin{equation} \sigma_{\min}(\mathbf{\Lambda}^{-1}) = \frac{1}{\lambda_N} \propto \frac{1}{N^2}. \end{equation}
    
    \item \textbf{The Frame Matrix (Exponential Scaling):} The term $\norm{\mat{W}_N^{-1}}_2$ is the reciprocal of the smallest singular value of the synthesis matrix $\mat{W}_N$, i.e., $\norm{\mat{W}_N^{-1}}_2 = 1/\sigma_{\min}(\mat{W}_N)$.
    It is a well-established result in the theory of Fourier extensions (see \cite[Sec.~4]{adcock2014numerical} and \cite{huybrechs2010fourier}) that the singular values of the synthesis operator for a Fourier extension frame decay exponentially to zero. For the discrete matrix $\mat{W}_N$ in the square case ($M=N$), this implies:
    \begin{equation} \sigma_{\min}(\mat{W}_N) \le C_1 e^{-\alpha(\delta, L)\, N}, \end{equation}
    where the decay rate $\alpha(\delta, L)$ depends on the extension ratio $\rho = L/(L+2\delta)$; specifically, $\alpha$ is related to $\log(1/\rho)$ and quantifies the degree of linear dependence among the restricted basis functions \cite[Sec.~4]{adcock2014numerical}. As $\delta$ increases, $\rho$ decreases, and the frame becomes more redundant, so $\alpha$ increases. Consequently, the norm of the inverse grows exponentially:
    \begin{equation} \norm{\mat{W}_N^{-1}}_2 \ge \frac{1}{C_1} e^{\alpha(\delta, L)\, N}. \end{equation}
\end{enumerate}

Combining these estimates, we obtain:
\begin{equation}
\norm{\mat{A}_P^{-1}}_2 \ge \frac{1}{C_2 N^2} e^{\alpha(\delta, L)\, N}.
\end{equation}
Since the exponential growth $e^{\alpha(\delta, L)\, N}$ dominates the polynomial decay $N^{-2}$, the condition number grows exponentially.
\end{proof}

\begin{remark}[The Implication of Square Systems]
The exponential growth of $\kappa(\mat{A}_P)$ confirms that for $M=N$, the method is numerically unstable; entries in the coefficient vector $\vect{c}$ can become arbitrarily large to represent $O(1)$ functions, leading to significant cancellation errors. While this can be mitigated by oversampling ($M > N$) to control $\sigma_{\min}(\mat{W}_N)$, the analysis of the square system reveals the ``native'' behavior of the operator-basis interaction. This baseline instability highlights the results of the following section, where the convection-diffusion operator exhibits a completely different structural behavior.
\end{remark}

%================================================================================
\section{Pre-Asymptotic Stability of the Convection-Diffusion Operator}
%================================================================================

Section 2 established that the extended-domain method is asymptotically unstable for the Poisson equation. Since the second-order diffusion term dominates the first-order convection term for high-frequency modes, the convection-diffusion operator inherits this same asymptotic exponential instability. This leaves a paradox: why is the method observed to be so much more stable in practice for convection-dominated problems?

This section resolves this paradox. We prove that the observed stability is a genuine, but \textbf{pre-asymptotic}, phenomenon. The mechanism is a fundamental structural difference in the discrete Green's functions. For convection-dominated problems, the inverse of the non-self-adjoint operator is \textbf{numerically highly asymmetric}, exhibiting exponential decay strictly in the upstream direction. This is in contrast to the \textbf{numerically dense} inverse of the self-adjoint Poisson operator. For any finite $N$, this structural property leads to a much smaller matrix norm and, consequently, a better-conditioned system.

We state our main theorem describing this structural dichotomy. To analyze convergence and uniform bounds as $N \to \infty$, the following proofs assume the method is stabilized (e.g., via oversampling $M \ge \beta N$ for a suitable constant $\beta > 1$, or via discrete regularization) so that the Lebesgue constant of the interpolation operator satisfies $\Lambda_N \le C$ uniformly in $N$. The generalization to the square system, where $\Lambda_N$ is no longer uniformly bounded, is given in Theorem~\ref{thm:square_system_decay} below.

\begin{theorem}[Upstream Exponential Decay in the Discrete Green's Function]
\label{thm:exp_decay_cd}
Let $\mat{L}_{CD,N} = \mat{A}_{CD,N}\mat{W}_N^{-1}$ be the physical-space representation of the constant-coefficient convection-diffusion ($k > 0$) operator on an ordered grid of collocation points $\{x_i\}$. Assume the discretization is stabilized such that $\Lambda_N \le C$ uniformly in $N$. Then for any $0 < \alpha < k$, there exists a constant $C_{\alpha} > 0$, independent of $N$, such that the entries of the inverse matrix $\mat{L}_{CD,N}^{-1}$ (the discrete Green's function) satisfy:
\begin{equation}|(\mat{L}_{CD,N}^{-1})_{ij}| \le C_{\alpha}\, \exp\bigl(-\alpha \max(0, x_j - x_i)\bigr).\end{equation}
The exponential bound is meaningful strictly in the upstream direction ($x_i \le x_j$), capturing the physical decay of the continuous Green's function. In the downstream direction ($x_i > x_j$), the flow carries information without attenuation, yielding an $O(1)$ dense block. The generalization to the square system ($M=N$), where $\Lambda_N$ grows exponentially, is given in Theorem~\ref{thm:square_system_decay}.
\end{theorem}

To prove Theorem \ref{thm:exp_decay_cd}, we proceed in three steps: we first establish the analytic decay of the continuous Green's function (Section 3.1), then prove that the collocation scheme is uniformly stable (Section 3.2, Lemmas \ref{lem:stability_proof}--\ref{lem:L_inv_bound}), and finally derive the entry-wise decay via a Combes--Thomas argument (Lemma \ref{lem:combes_thomas}).

%--------------------------------------------------------------------------------
\subsection{The Continuous Green's Function}
%--------------------------------------------------------------------------------

The physical-space inverse matrix $\mat{L}_N^{-1}$ is the exact discrete equivalent of the continuous Green's function $G(x,y)$. The structure of $\mat{L}_N^{-1}$ is therefore dictated by the underlying physics of the differential operator $\mathcal{L}$.

For the 1D Poisson operator $\mathcal{L}_P = -\Delta$ on a domain $I = [0, L]$ with homogeneous Dirichlet boundaries, the continuous Green's function is a global, symmetric, piecewise linear function:
\begin{equation}G_P(x,y) = \begin{cases} x(L-y)/L, & x \le y \\ y(L-x)/L, & x > y \end{cases}\end{equation}
This function does not exhibit exponential decay. Consequently, the discrete Poisson Green's function $\mat{L}_{P,N}^{-1}$ is a dense matrix.

In contrast, for the convection-diffusion operator $\mathcal{L}_{CD} = -\Delta + k \frac{d}{dx}$ (assuming $k > 0$ for rightward flow), the continuous Green's function $G_{CD}(x,y)$ is highly asymmetric. Due to the first-order convective derivative, the upstream tail of the Green's function (where $x < y$) decays exactly at an exponential rate. Specifically, the homogeneous solution involves the roots of the characteristic equation $\lambda^2 - k\lambda = 0$, leading to terms like $e^{kx}$. The upstream influence decays proportional to $e^{-k|x-y|}$.

%--------------------------------------------------------------------------------
\subsection{Discrete Inheritance via Stability}
%--------------------------------------------------------------------------------

Standard spectral convergence theory dictates that if a numerical discretization is stable, the discrete Green's function $\mat{L}_N^{-1}$ converges pointwise to the continuous Green's function $G(x,y)$ as $N \to \infty$. Therefore, to prove that $\mat{L}_{CD,N}^{-1}$ exhibits exponential decay (Theorem \ref{thm:exp_decay_cd}), it is sufficient to prove that the collocation scheme is uniformly stable in the maximum norm.

\begin{lemma}[Stability of the Stabilized Collocation Scheme]
\label{lem:stability_proof}
Let $u_N \in V_N$ be the solution of the extended-domain spectral collocation method for the operator $\mathcal{L}_{CD}$ with forcing data given by the nodal values of an interpolant $f_I$. Assuming the discretization is sufficiently stabilized (e.g., via oversampling $M \ge \beta N$) such that the quadrature aliasing error is controlled, the scheme is uniformly stable in the sense that there exists a constant $C_{num}$, independent of $N$, such that
\begin{equation} \norm{u_N}_{H^1(I)} \le C_{num} \norm{f_I}_{L^2(I)}. \end{equation}
\end{lemma}
\begin{proof}
The proof uses a contradiction argument, correctly founded on a discrete G\r{a}rding inequality.

\noindent\textbf{1. Discrete G\r{a}rding Inequality.}
Let $a_Q(u,v) = Q((\mathcal{L}_{CD}u)\bar{v})$ be the discrete bilinear form, where $Q(\cdot)$ is the quadrature rule. Let $a(u,v)$ be the corresponding continuous form. The continuous operator satisfies a G\r{a}rding inequality: $\text{Re}[a(u_N, u_N)] \ge \alpha_0 |u_N|_{H^1}^2 - C_0 \norm{u_N}_{L^2}^2$. We show the discrete form inherits this property by bounding the quadrature error, $E_N(g) = Q(g) - I(g)$, where $g = (\mathcal{L}_{CD}u_N)\bar{u}_N$.
\begin{equation} \text{Re}[a_Q(u_N, u_N)] = \text{Re}[a(u_N, u_N)] + \text{Re}[E_N(g)]. \end{equation}
The integrand $g$ is a trigonometric polynomial of degree at most $2N$, and the error $E_N(g)$ arises from aliasing. The key to bounding this error is to note that it can be controlled by the norms of the functions being multiplied. A standard result from the analysis of pseudospectral methods, derived by decomposing the solution into low- and high-frequency components, provides the following bound (see the argument in \cite[Sec.~6.4]{canuto2006spectral}). For any $\epsilon > 0$, there exists a constant $C_\epsilon$, independent of $N$, such that:
\begin{equation} |E_N((\mathcal{L}_{CD}u_N)\bar{u}_N)| \le \epsilon |u_N|_{H^1}^2 + C_\epsilon \norm{u_N}_{L^2}^2. \end{equation}
This inequality holds because the aliasing error can be made arbitrarily small in the highest-order norm ($|u_N|_{H^1}^2$) at the expense of a larger factor in the lower-order norm ($\norm{u_N}_{L^2}^2$), with $\epsilon$ controlling the trade-off. By substituting this bound into our expression for the discrete form, we get:
\begin{equation} \text{Re}[a_Q(u_N, u_N)] \ge (\alpha_0 |u_N|_{H^1}^2 - C_0 \norm{u_N}_{L^2}^2) - (\epsilon |u_N|_{H^1}^2 + C_\epsilon \norm{u_N}_{L^2}^2). \end{equation}
Rearranging the terms yields:
\begin{equation} \text{Re}[a_Q(u_N, u_N)] \ge (\alpha_0 - \epsilon)|u_N|_{H^1}^2 - (C_0 + C_\epsilon)\norm{u_N}_{L^2}^2. \end{equation}
We can choose a fixed, small value for $\epsilon$, for instance $\epsilon = \alpha_0/2$. This choice fixes the associated constant $C_\epsilon$. We then arrive at the discrete G\r{a}rding inequality with constants $\alpha_1 = \alpha_0/2$ and $C_1 = C_0+C_{\alpha_0/2}$, which are independent of $N$:
\begin{equation} \text{Re}[a_Q(u_N, u_N)] \ge \alpha_1 |u_N|_{H^1}^2 - C_1 \norm{u_N}_{L^2}^2. \end{equation}

\noindent\textbf{2. Proof by Contradiction.} Assume the stability estimate is false. Then there exists a sequence $\{u_N\}_{N=1}^\infty$ with $u_N \in V_N$ such that $\norm{u_N}_{H^1(I)} = 1$ for all $N$, while the corresponding forcing terms $f_{I,N} = \mathcal{L}_{CD}u_N$ satisfy $\norm{f_{I,N}}_{L^2(I)} \to 0$ as $N \to \infty$.

From the G\r{a}rding inequality and the scheme $a_Q(u_N, u_N) = (f_{I,N}, u_N)_Q$, we have
\begin{equation} \alpha_1 |u_N|_{H^1}^2 - C_1 \norm{u_N}_{L^2}^2 \le |(f_{I,N}, u_N)_Q| \le \norm{f_{I,N}}_Q \norm{u_N}_Q \le C_Q^2 \norm{f_{I,N}}_{L^2} \norm{u_N}_{L^2}, \end{equation}
where we have invoked Cauchy-Schwarz and the equivalence of the discrete and continuous norms on $V_N$ with constant $C_Q$. Since $\norm{u_N}_{H^1}=1$, its norm components $|u_N|_{H^1}$ and $\norm{u_N}_{L^2}$ are bounded. As $N\to\infty$, the RHS tends to zero. Thus, $\limsup_{N\to\infty} (\alpha_1 |u_N|_{H^1}^2 - C_1 \norm{u_N}_{L^2}^2) \le 0$.
We use the identity $|u_N|_{H^1}^2 = \norm{u_N}_{H^1}^2 - \norm{u_N}_{L^2}^2 = 1 - \norm{u_N}_{L^2}^2$. Substituting this into the inequality gives:
\begin{equation} \alpha_1 (1 - \norm{u_N}_{L^2}^2) - C_1 \norm{u_N}_{L^2}^2 \le o(1), \end{equation}
where $o(1)$ represents a term that vanishes as $N\to\infty$. Rearranging the terms, we find
\begin{equation} \alpha_1 - o(1) \le (\alpha_1 + C_1) \norm{u_N}_{L^2}^2. \end{equation}
This implies $\liminf_{N\to\infty} \norm{u_N}_{L^2}^2 \ge \frac{\alpha_1}{\alpha_1 + C_1} > 0$. This step is valid regardless of the relative sizes of $\alpha_1$ and $C_1$.

\noindent\textbf{3. The Contradiction.} The sequence $\{u_N\}$ is bounded in $H^1(I)$. By the Rellich-Kondrachov theorem, a subsequence (still denoted $\{u_N\}$) exists that converges strongly in $L^2(I)$ to a limit $u$, and weakly in $H^1(I)$ to $u$. Since $\liminf \norm{u_N}_{L^2} > 0$, the limit function is non-zero, $\norm{u}_{L^2} \ne 0$.

We now show that $u$ is a weak solution to $\mathcal{L}_{CD}u=0$. Fix any smooth test function $v \in C_c^\infty(I)$. We decompose:
\begin{equation} a_Q(u_N, \mathcal{I}_N v) - a(u, v) = \underbrace{[a_Q(u_N, \mathcal{I}_N v) - a(u_N, \mathcal{I}_N v)]}_{E_Q} + \underbrace{[a(u_N, \mathcal{I}_N v) - a(u_N, v)]}_{E_I} + \underbrace{[a(u_N, v) - a(u, v)]}_{E_W}. \end{equation}
\textbf{Term $E_W$:} Since $a(\cdot, v)$ is a bounded linear functional on $H^1(I)$ (because $v$ is fixed and smooth), weak convergence $u_N \rightharpoonup u$ in $H^1$ gives $E_W \to 0$.

\textbf{Term $E_I$:} Because $v \in C_c^\infty(I)$ is fixed, $\mathcal{I}_N v \to v$ in $H^1(I)$ at spectral rate. The bilinear form satisfies $|a(u_N, \mathcal{I}_N v - v)| \le C_a \norm{u_N}_{H^1} \norm{\mathcal{I}_N v - v}_{H^1} \le C_a \norm{\mathcal{I}_N v - v}_{H^1} \to 0$.

\textbf{Term $E_Q$:} This is the quadrature error $Q((\mathcal{L}_{CD}u_N)\overline{\mathcal{I}_N v}) - \int (\mathcal{L}_{CD}u_N)\overline{\mathcal{I}_N v}\,dx$. We decompose $u_N = u_N^{\le} + u_N^{>}$ into low- and high-frequency modes. For a fixed smooth test function $v$, $\mathcal{I}_N v$ is uniformly bounded and its high-frequency spectrum decays super-exponentially. For the low-frequency part $u_N^{\le}$, the integrand is a smooth, band-limited function; thus, the oversampled $M$-point quadrature integrates it with super-exponential accuracy (error $O(e^{-cN}) \to 0$). For the high-frequency part, we invoke the $H^1$-boundedness of $u_N$ and the rapid spectral decay of $\mathcal{I}_N v$ to bound the remaining aliasing contributions by $O(N^{-1}) \to 0$ (cf.\ \cite[Sec.~6.4]{canuto2006spectral}).

Since $a_Q(u_N, \mathcal{I}_N v) = (f_{I,N}, \mathcal{I}_N v)_Q \to 0$ (because $\norm{f_{I,N}}_{L^2} \to 0$ and $\norm{\mathcal{I}_N v}_{L^2}$ is bounded), we conclude $a(u,v) = 0$ for all $v \in C_c^\infty(I)$. This means $\mathcal{L}_{CD}u=0$ weakly. However, the continuous problem is well-posed, with $u=0$ as its only solution. This contradicts $\norm{u}_{L^2} \ne 0$. Therefore, the scheme is stable.
\end{proof}

\begin{remark}[Failure of Uniform G\r{a}rding for the Square System]
\label{rem:H1_square}
The uniform $H^1$ stability established above critically relies on bounding the quadrature error $E_N(g)$. For the square system ($M=N$), the $N$-point collocation grid cannot resolve the maximum frequency ($2N$) of the product $g = (\mathcal{L}_{CD}u_N)\bar{u}_N$. Consequently, the aliasing error is $O(1)$ and cannot be bounded by a small fraction $\epsilon$ of the $H^1$ norm. The discrete G\r{a}rding inequality therefore fails for the unstabilized square system, and uniform stability requires stabilization (e.g., $M > N$). However, we can still establish a weaker analytical bound for the square system by bypassing the G\r{a}rding inequality entirely, as shown in Theorem~\ref{thm:square_system_decay}.
\end{remark}

\begin{lemma}[Boundedness of the Inverse Physical-Space Operator]
\label{lem:L_inv_bound}
For the stabilized convection-diffusion operator $\mathcal{L}_{CD} = -d^2/dx^2 + k(x)d/dx$ with a bounded and strictly positive convection coefficient $k(x)$, the inverse physical-space operators satisfy
\begin{equation}\norm{\mat{L}_{CD,N}^{-1}}_\infty \le C_{inv},\end{equation}
where $C_{inv}$ depends on $C_{emb}$, $C_{num}$, and $L$, but is independent of $N$.
\end{lemma}
\begin{proof}
The proof relies on the stability of the numerical scheme established in Lemma \ref{lem:stability_proof}. Let $\vect{u} = \mat{L}_{CD,N}^{-1} \vect{f}$, which corresponds to the solution $u_N \in V_N$ for an interpolated forcing function $f_I(x)$. From Lemma \ref{lem:stability_proof}, we have the stability estimate
\begin{equation}
\norm{u_N}_{H^1(I)} \le C_{num} \norm{f_I}_{L^2(I)}.
\end{equation}
The proof proceeds by connecting the discrete $\ell_\infty$-norm to the continuous norms. By the one-dimensional Sobolev embedding theorem, $\norm{u_N}_{L^\infty(I)} \le C_{emb} \norm{u_N}_{H^1(I)}$. The solution vector is bounded by $\norm{\vect{u}}_\infty \le \norm{u_N}_{L^\infty(I)}$. Since the scheme is stabilized, the Lebesgue constant is uniformly bounded ($\Lambda_N \le C$). Thus, the forcing term is bounded by: $\norm{f_I}_{L^2(I)} \le \sqrt{L}\norm{f_I}_{L^\infty(I)} \le \sqrt{L} C \norm{\vect{f}}_\infty$.
\begin{equation}
\norm{\vect{u}}_\infty \le C_{emb} C_{num} \sqrt{L}\, C \norm{\vect{f}}_\infty.
\end{equation}
The induced $\ell_\infty$-norm of the inverse is therefore:
\begin{equation} \norm{\mat{L}_{CD,N}^{-1}}_\infty = \sup_{\vect{f} \neq 0} \frac{\norm{\vect{u}}_\infty}{\norm{\vect{f}}_\infty} \le C_{inv}, \end{equation}
where $C_{inv} = C_{emb} C_{num} \sqrt{L}\, C$ is independent of $N$.
\end{proof}

Lemmas \ref{lem:stability_proof} and \ref{lem:L_inv_bound} establish $\ell_\infty$ boundedness of the inverse operator, with the Lebesgue constant $\Lambda_N$ appearing explicitly as a prefactor. To complete the proof of Theorem \ref{thm:exp_decay_cd}, we must bridge this operator-norm bound to the \emph{entry-wise} exponential decay. The following lemma achieves this via a Combes--Thomas conjugation argument.

\begin{lemma}[Entry-Wise Asymmetric Decay via Weighted Stability]
\label{lem:combes_thomas}
Assume the discretization is stabilized such that the Lebesgue constant is uniformly bounded ($\Lambda_N \le C$). For any $0 < \alpha < |k|$, there exist constants $C_{\alpha} > 0$ and $N_0 \in \mathbb{N}$, independent of $N$, such that for all $N \ge N_0$:
\begin{equation}|(\mat{L}_{CD,N}^{-1})_{ij}| \le C_{\alpha}\, \exp\bigl(-\alpha \max(0, x_j - x_i)\bigr).\end{equation}
\end{lemma}
\begin{proof}
The proof adapts the classical Combes--Thomas argument \cite{combes1973asymptotic} to the spectral collocation setting. Fix $0 < \alpha < |k|$ and a column index $j$.

\noindent\textbf{1. Conjugated Operator.}
For the upstream bound ($x_i < x_j$, assuming $k > 0$), define the smooth weight function $\mu(x) = e^{-\alpha(x - x_j)}$ and the corresponding diagonal matrix $\mat{D}_{-\alpha}$ with $(\mat{D}_{-\alpha})_{ii} = e^{-\alpha(x_i - x_j)}$. Consider the conjugated operator:
\begin{equation} \widetilde{\mat{L}} = \mat{D}_{-\alpha} \mat{L}_{CD,N} \mat{D}_{-\alpha}^{-1}. \end{equation}
At the continuous level, conjugating $\mathcal{L}_{CD} = -d^2/dx^2 + k\,d/dx$ by $e^{-\alpha(x - x_j)}$ transforms it into a modified operator. If $u = e^{\alpha(x-x_j)}\tilde{u}$, then direct computation gives:
\begin{equation} \mathcal{L}_{CD}[e^{\alpha(x-x_j)}\tilde{u}] = e^{\alpha(x-x_j)}\bigl[-\tilde{u}'' + (k - 2\alpha)\tilde{u}' + \alpha(k - \alpha)\tilde{u}\bigr]. \end{equation}
The conjugated operator is therefore $\widetilde{\mathcal{L}} = -d^2/dx^2 + (k - 2\alpha)\,d/dx + \alpha(k - \alpha)$. Since $0 < \alpha < k$, the zeroth-order (reaction) coefficient $c_0 = \alpha(k - \alpha)$ is strictly positive.

\noindent\textbf{2. Improved Coercivity of the Conjugated Operator.}
The G\r{a}rding inequality for $\widetilde{\mathcal{L}}$ takes the form:
\begin{equation} \text{Re}[\tilde{a}(u,u)] \ge \tilde{\alpha}_0 |u|_{H^1}^2 + c_0 \norm{u}_{L^2}^2 - \tilde{C}_0 \norm{u}_{L^2}^2 = |u|_{H^1}^2 + (c_0 - \tilde{C}_0)\norm{u}_{L^2}^2, \end{equation}
where $\tilde{\alpha}_0 = 1$ (from the diffusion term) and $\tilde{C}_0 = \frac{1}{2}\norm{(k - 2\alpha)'}_{L^\infty} = 0$ for constant $k$. Thus, with constant coefficients, $\tilde{C}_0 = 0$ and $c_0 > 0$, so the conjugated operator satisfies a G\r{a}rding inequality with the strictly positive lower bound:
\begin{equation} \text{Re}[\tilde{a}(u,u)] \ge |u|_{H^1}^2 + \alpha(k-\alpha)\norm{u}_{L^2}^2 \ge \min(1, \alpha(k-\alpha))\norm{u}_{H^1}^2. \end{equation}
This is a genuine coercivity estimate (not merely a G\r{a}rding inequality), which strengthens the stability analysis.

\noindent\textbf{3. Stability of the Conjugated Scheme.}
Note that the algebraic conjugation $\widetilde{\mat{L}}$ differs from the exact spectral collocation matrix of $\widetilde{\mathcal{L}}$ due to the commutation error between spectral differentiation and multiplication by the weight $e^{\pm \alpha x}$. Let $\mat{M}_N$ denote the exact spectral collocation matrix of the continuous conjugated operator $\widetilde{\mathcal{L}}$. Because the exponential weight function is entire, the interpolation error of the weighted function $e^{\alpha(x-x_j)} \tilde{u}_N(x)$ converges to zero super-exponentially. Consequently, the matrix difference $\widetilde{\mat{L}} - \mat{M}_N$ is bounded by an error matrix $\mat{E}_{comm}$ whose norm decays to zero as $N \to \infty$. Since $\mat{M}_N$ inherits the improved coercivity of the continuous operator $\widetilde{\mathcal{L}}$ (by the same argument as Lemma~\ref{lem:stability_proof}), and $\widetilde{\mat{L}} = \mat{M}_N + \mat{E}_{comm}$ is a small perturbation, $\widetilde{\mat{L}}$ is invertible for sufficiently large $N$ with $\norm{\widetilde{\mat{L}}^{-1}}_\infty \le \tilde{C}_{inv}$, where $\tilde{C}_{inv}$ depends on $\alpha$ and $k$ but is independent of $N$.

\noindent\textbf{4. Entry-Wise Bound.}
Since $\mat{L}_{CD,N}^{-1} = \mat{D}_{-\alpha}^{-1} \widetilde{\mat{L}}^{-1} \mat{D}_{-\alpha}$, the entry $(i,j)$ satisfies:
\begin{equation} (\mat{L}_{CD,N}^{-1})_{ij} = e^{\alpha(x_i - x_j)} (\widetilde{\mat{L}}^{-1} \mat{D}_{-\alpha} \mat{e}_j)_i = e^{\alpha(x_i - x_j)} (\widetilde{\mat{L}}^{-1} \mat{e}_j)_i, \end{equation}
where the last equality uses $(\mat{D}_{-\alpha})_{jj} = 1$. From Step 3, we have the uniform bound $\norm{\widetilde{\mat{L}}^{-1}}_\infty \le \tilde{C}_{inv}$. This implies that the elements of the conjugated inverse matrix are globally bounded by $|(\widetilde{\mat{L}}^{-1})_{ij}| \le \tilde{C}_{inv}$.

Finally, we translate this property of the conjugated matrix back to the original matrix $\mat{L}_{CD,N}^{-1}$. Recalling the relation established earlier, $(\mat{L}_{CD,N}^{-1})_{ij} = e^{\alpha(x_i - x_j)} (\widetilde{\mat{L}}^{-1})_{ij}$, we take the absolute value and apply the uniform bound for the conjugated matrix entries. This directly yields the exponential decay for the upstream direction ($x_i \le x_j$, so $x_i-x_j \le 0$):
\begin{equation}
|(\mat{L}_{CD,N}^{-1})_{ij}| \le \tilde{C}_{inv} e^{\alpha (x_i-x_j)}.
\end{equation}
For the downstream direction ($x_i > x_j$), the conjugated operator bound yields $\tilde{C}_{inv} e^{\alpha(x_i - x_j)}$, which grows exponentially. Instead, we simply rely on the unweighted stability bound from Lemma~\ref{lem:L_inv_bound}, which guarantees $|(\mat{L}_{CD,N}^{-1})_{ij}| \le \norm{\mat{L}_{CD,N}^{-1}}_\infty \le C_{inv}$. 
This mathematically confirms that the discrete Green's function does not decay downstream, but rather plateaus to an $O(1)$ constant, consistent with the physics of the continuous convection-diffusion operator where information is carried downstream without attenuation.
Combining these two regimes yields the asymmetric bound $|(\mat{L}_{CD,N}^{-1})_{ij}| \le C_{\alpha}\, \exp(-\alpha \max(0, x_j - x_i))$ with $C_\alpha = \max(\tilde{C}_{inv}, C_{inv})$.
\end{proof}

\begin{remark}[Scope of Lemmas~\ref{lem:stability_proof}--\ref{lem:combes_thomas}]
\label{rem:variable_k}
While Theorem~\ref{thm:exp_decay_cd} is stated for constant $k$ (where the continuous Green's function is explicit and the Combes--Thomas conjugation yields exact formulas), the stability estimates in Lemmas~\ref{lem:stability_proof} and \ref{lem:L_inv_bound} hold for variable-coefficient operators $\mathcal{L}_{CD} = -d^2/dx^2 + k(x)d/dx$ with $k(x) \ge k_{\min} > 0$. For variable coefficients, the Combes--Thomas argument in Lemma~\ref{lem:combes_thomas} produces a modified operator with spatially varying reaction coefficient $\alpha(k(x) - \alpha)$, which remains positive when $\alpha < k_{\min}$. The exponential decay rate is then governed by $k_{\min}$ rather than $k$.
\end{remark}

\begin{theorem}[Row-Sum Bound for the Square System]
\label{thm:square_system_decay}
Let $\mat{L}_N$ be the physical-space operator matrix for the square system ($M=N$) for either the Poisson ($\mathcal{L}_P$) or convection-diffusion ($\mathcal{L}_{CD}$) operator. Its maximum row sum (the $\ell_\infty$-norm) is bounded by the test-space Lebesgue constant $\Lambda_N$:
\begin{equation} \norm{\mat{L}_N^{-1}}_\infty \le C \Lambda_N, \end{equation}
where $C$ depends only on the continuous inverse operator.
\end{theorem}
\begin{proof}
Assume boundary conditions are strongly enforced via basis recombination, so the trial space $\Phi_N \subset V_N$ consists entirely of functions satisfying homogeneous Dirichlet boundaries. Since the continuous operator $\mathcal{L}$ with homogeneous Dirichlet conditions is injective (the only solution to $\mathcal{L}u = 0$ with $u(0)=u(L)=0$ is $u=0$, by the maximum principle), the map $\mathcal{L}: \Phi_N \to \Psi_N := \mathcal{L}(\Phi_N)$ is a bijection, and $\mathcal{L}^{-1}$ maps $\Psi_N$ back to $\Phi_N$. Let $\mathcal{P}_N$ be the exact interpolation operator at the interior collocation points into the image space $\Psi_N$. For the square system, collocation enforces exactly $\mathcal{L} u_N = \mathcal{P}_N f$ because both sides belong to $\Psi_N$ and match at the collocation points. Because the continuous inverse $\mathcal{L}^{-1}: L^\infty(I) \to L^\infty(I)$ is a bounded operator (with norm $C = \norm{\mathcal{L}^{-1}}_\infty$), the continuous solution satisfies $u_N(x) = \mathcal{L}^{-1}(\mathcal{P}_N f)(x)$.
Evaluating this at the collocation nodes gives the discrete physical-space solution vector $\vect{u}$:
\begin{equation} \norm{\vect{u}}_\infty \le \norm{u_N}_{L^\infty(I)} \le \norm{\mathcal{L}^{-1}}_\infty \norm{\mathcal{P}_N f}_{L^\infty(I)}. \end{equation}
The interpolation operator $\mathcal{P}_N$ maps into $\Psi_N$, whose Lebesgue constant $\tilde{\Lambda}_N$ is asymptotically equivalent to $\Lambda_N$ (as noted in Definition 3, since the second-derivative term dominates for high frequencies). Thus, $\norm{\mathcal{P}_N f}_{L^\infty(I)} \le \tilde{\Lambda}_N \norm{\vect{f}}_\infty$, and the discrete inverse matrix satisfies $\norm{\mat{L}_N^{-1}}_\infty = \sup (\norm{\vect{u}}_\infty / \norm{\vect{f}}_\infty) \le C \tilde{\Lambda}_N = O(\Lambda_N)$.
\end{proof}

While the discrete $M=N$ system lacks the coercivity needed for sharp entry-wise bounds, the structural dichotomy between the operators can be identified by analyzing how their continuous Green's functions propagate boundary-localized errors, which are characteristic of the ill-conditioned Fourier extension frame.

\begin{theorem}[Continuous Attenuation of Boundary Instabilities]
\label{thm:boundary_attenuation}
Let $e_{bnd}(x)$ be an error mode supported in $[L-\epsilon, L]$ with mass $\norm{e_{bnd}}_{L^1}$. 
For the Poisson operator, the continuous integral $\mathcal{L}_P^{-1} e_{bnd}(x)$ propagates into the interior domain $x < L-\epsilon$, bounded only algebraically as $|\mathcal{L}_P^{-1} e_{bnd}(x)| \le \frac{x\epsilon}{L} \norm{e_{bnd}}_{L^1}$.
For the convection-diffusion operator with constant $k > 0$, the continuous inverse strictly attenuates the error:
\begin{equation} |\mathcal{L}_{CD}^{-1} e_{bnd}(x)| \le \frac{e^{k\epsilon}}{k(1-e^{-kL})} e^{-k(L-x)} \norm{e_{bnd}}_{L^1}. \end{equation}
Consequently, the error in the interior decays exponentially away from the boundary.
\end{theorem}
\begin{proof}
For a point $x < L-\epsilon$ in the interior, we evaluate the action of the continuous inverse operators via integration against their Green's functions for $y > x$.

For the Poisson operator, the continuous Green's function for $y \ge x$ is $G_P(x,y) = \frac{x(L-y)}{L}$. Applying this to the boundary error yields:
\begin{equation} |\mathcal{L}_P^{-1} e_{bnd}(x)| \le \int_{L-\epsilon}^L G_P(x,y) |e_{bnd}(y)| \,dy \le \left(\max_{y \in [L-\epsilon, L]} \frac{x(L-y)}{L}\right) \norm{e_{bnd}}_{L^1} = \frac{x\epsilon}{L} \norm{e_{bnd}}_{L^1}. \end{equation}
This bound depends algebraically on $\epsilon$ and $x$, meaning the error spreads globally into the interior without spatial attenuation.

For the convection-diffusion operator with $k > 0$, the continuous Green's function for $y \ge x$ is:
\begin{equation} G_{CD}(x,y) = \frac{1 - e^{-kx}}{k(1 - e^{-kL})} \left( e^{-k(y-x)} - e^{-k(L-x)} \right). \end{equation}
Since $1 - e^{-kx} < 1$ and $e^{-k(L-x)} > 0$, this is strictly bounded from above:
\begin{equation} G_{CD}(x,y) \le \frac{1}{k(1 - e^{-kL})} e^{-k(y-x)}. \end{equation}
Integrating against $e_{bnd}(y)$ supported in $[L-\epsilon, L]$ yields:
\begin{equation} |\mathcal{L}_{CD}^{-1} e_{bnd}(x)| \le \int_{L-\epsilon}^L G_{CD}(x,y) |e_{bnd}(y)| \,dy \le \left( \max_{y \in [L-\epsilon, L]} G_{CD}(x,y) \right) \norm{e_{bnd}}_{L^1}. \end{equation}
For fixed $x$, the maximum over $y \in [L-\epsilon, L]$ occurs at $y = L-\epsilon$. Thus:
\begin{equation} |\mathcal{L}_{CD}^{-1} e_{bnd}(x)| \le \frac{e^{-k(L-\epsilon-x)}}{k(1 - e^{-kL})} \norm{e_{bnd}}_{L^1} = \frac{e^{k\epsilon}}{k(1 - e^{-kL})} e^{-k(L-x)} \norm{e_{bnd}}_{L^1}. \end{equation}
As the observation point $x$ moves into the interior, the boundary error is exponentially attenuated by $e^{-k(L-x)}$, effectively shielding the interior from the downstream frame instability.
\end{proof}

\begin{corollary}[Explicit Pre-Asymptotic Regime]
\label{cor:pre_asymptotic_threshold}
Combining the spatial attenuation of Theorem~\ref{thm:boundary_attenuation} with the frame growth $O(e^{c(\delta,L)N})$, the method exhibits a pre-asymptotic stable regime when the physical decay limits the numerical instability. Errors from the downstream boundary remain small in the interior whenever $k(L-x) > c(\delta, L) N$.
This yields a critical resolution $N^* \sim k L / c(\delta, L)$, below which the interior solution retains accuracy. For the Poisson operator ($k=0$), there is no exponential attenuation, so $N^* = 0$, and the method is unconditionally unstable.
\end{corollary}

\begin{figure}[ht!]
    \centering
    \includegraphics[width=0.9\textwidth]{matrix_structures.png}
    \caption{Direct visualization of the structural dichotomy for $N=128$, $\delta=0.05$, and $k_0=50$. The colormaps show $\log_{10}$ of the absolute value of the matrix entries. (a), (c): Both forward operators $\mat{L}_{P}$ and $\mat{L}_{CD}$ are localized. (b): The Poisson Green's function $G_P$ is dense. (d): The convection-diffusion Green's function $G_{CD}$ is numerically highly asymmetric. (e): Self-normalized cross-section of row $i=64$: each curve plots $|G_{ij}|$ versus distance $|j-i|$. The Poisson entries (blue) decay only algebraically. The C-D upstream influence (red solid) drops exponentially, while the downstream influence (orange dashed) plateaus to the smooth $O(1/k)$ physical limit, demonstrating the fundamental asymmetry of the convection-diffusion Green's function.}
    \label{fig:matrix_structures}
\end{figure}

\begin{remark}[Visual Verification of the Green's Function Structure]
The structural dichotomy is visible in numerical computations. Figure \ref{fig:matrix_structures} displays the magnitude of the entries for the physical-space operators and their Green's functions.
\begin{itemize}
    \item \textbf{Poisson Green's Function (Fig.\ \ref{fig:matrix_structures}b, e):} The Green's function $G_P$ is dense: after self-normalization, its entries decay only about two orders of magnitude over $|j-i| \in [1, N/2]$. This slow decay reflects the global nature of the elliptic Green's function.
    \item \textbf{Convection-Diffusion Green's Function (Fig.\ \ref{fig:matrix_structures}d, e):} In contrast, $G_{CD}$ is highly asymmetric. Panel (e) separates the upstream (red solid) and downstream (orange dashed) influences. The upstream influence exhibits true exponential decay to zero. The downstream influence drops initially from the numerical noise peak but quickly plateaus to a smooth $O(1)$ constant, consistent with the continuous physics.
\end{itemize}
This structural asymmetry in the inverse operator is the physical mechanism that suppresses the accumulation of round-off errors and stabilizes the method in the pre-asymptotic regime, distinct from the behavior of the underlying ill-conditioned basis.
\end{remark}

\begin{remark}[Pre-Asymptotic Conditioning and the Structural Dichotomy]
\label{rem:pre_asymptotic}
The distinct practical stability between the Poisson and convection-diffusion operators lies in their structural asymmetry, not their row sums.
By Theorem~\ref{thm:square_system_decay}, both discrete physical operators have row sums bounded by $O(\Lambda_N)$. The $\ell_\infty$-norm of the full collocation system $\mat{A}^{-1} = \mat{W}_N^{-1}\mat{L}_N^{-1}$ satisfies the sub-multiplicative bound $\norm{\mat{A}^{-1}}_\infty \le \norm{\mat{W}_N^{-1}}_\infty \norm{\mat{L}_N^{-1}}_\infty$. 

However, this crude bound is pessimistic for convection-diffusion. The synthesis inverse $\mat{W}_N^{-1}$ is exponentially large and dense, mapping boundary instabilities globally. For the Poisson equation, the discrete Green's function $\mat{L}_{P,N}^{-1}$ is structurally dense. When multiplied together, this dense matrix mixes and amplifies the exponentially growing singular modes of $\mat{W}_N^{-1}$, destabilizing the solution.

In contrast, numerical evidence shows the inverse of the convection-diffusion operator, $\mat{L}_{CD,N}^{-1}$, inherits the highly asymmetric structure of the continuous operator. When computing $\mat{W}_N^{-1}\mat{L}_{CD,N}^{-1}$, the instability is suppressed via a two-part mechanism:
\begin{enumerate}
    \item \textbf{Upstream Attenuation:} Boundary errors at the downstream boundary ($x=L$) are exponentially attenuated as they propagate upstream, shielding the interior.
    \item \textbf{Downstream Smoothing:} Boundary errors at the upstream boundary ($x=0$) propagate downstream. However, because the Green's function is a smooth constant downstream, it acts as an integrator. When integrated against the highly oscillatory frame errors (from $\mat{W}_N^{-1}$), the errors average out to near-zero.
\end{enumerate}

This structural combination of exponential attenuation and smooth integration filters the high-frequency noise inherent to the frame, granting the convection-diffusion system improved practical stability in the pre-asymptotic regime.
\end{remark}

\begin{remark}[Practical Implications for Stabilization]
This structural dichotomy has direct practical consequences for solver design. Because the highly asymmetric inverse of the convection-diffusion operator effectively filters the high-frequency noise inherent to the ill-conditioned Fourier extension frame, this operator requires significantly less oversampling or regularization parameter tuning than the Poisson operator to achieve stable, high-accuracy solutions in the pre-asymptotic regime.
\end{remark}

\begin{remark}[The P\'eclet Number and the Transition to Instability]
The transition from the stable pre-asymptotic regime to the true asymptotic instability is governed by the competition between convection and diffusion, which can be quantified by a \textbf{modal P\'eclet number}, $Pe_j \propto |k|/j$.
\begin{itemize}
    \item \textbf{Convection-Dominated Modes ($j \ll |k|$):} For low-frequency modes, the operator behaves like a first-order transport operator. The inverse matrix is dominated by these modes and exhibits the favorable highly asymmetric structure, leading to the observed stability.
    \item \textbf{Diffusion-Dominated Modes ($j \gg |k|$):} For high-frequency modes, the $j^2$ scaling of the diffusive part dominates the $j$ scaling of the convective part. These modes behave like those of the unstable Poisson problem.
\end{itemize}
For any fixed $N$, if $k$ is large enough, most modes are in the convection-dominated regime. However, as $N \to \infty$, the basis is inevitably populated by diffusion-dominated modes ($j \approx N$). It is these highest-frequency modes that ultimately drive the true exponential instability of the system, but their effect is only felt when $N$ becomes large enough to overcome the stabilizing influence of the lower-frequency part of the basis. This fully reconciles the theory with the numerical observations.
\end{remark}

%================================================================================
\section{Extension to Higher Dimensions}
%================================================================================

The structural dichotomy observed in one dimension extends to $d$-dimensional problems on tensor-product domains. The analysis here reveals a tension: while the baseline ill-conditioning of the linear system becomes more pronounced due to the tensor-product nature of the basis, the favorable pre-asymptotic properties of the convection-diffusion operator are strictly preserved.

%--------------------------------------------------------------------------------
\subsection{Framework in $d$-Dimensions}
%--------------------------------------------------------------------------------

We consider a physical domain $R_0 = (0, L_1) \times \dots \times (0, L_d)$ embedded within a larger computational domain $R = (-\delta_1, L_1+\delta_1) \times \dots \times (-\delta_d, L_d+\delta_d)$. The framework is built upon tensor products of the one-dimensional components.

\begin{itemize}
    \item \textbf{Basis Functions:} The basis for the space $V_{N,d}$ is formed by the tensor products of the 1D eigenfunctions, indexed by a multi-index $\mathbf{j}=(j_1, \dots, j_d)$, where $1 \le j_k \le N$ for each $k=1,\dots,d$:
    \begin{equation}w_{\mathbf{j}}(\mathbf{x}) = \prod_{k=1}^d w_{j_k}(x_k).\end{equation}
    There are $M=N^d$ such basis functions. They are the eigenfunctions of the negative Laplacian on the extended domain $R$ with homogeneous Dirichlet conditions.
    
    \item \textbf{Eigenvalues:} The corresponding eigenvalues of $-\Delta$ on $R$ are the sums of the 1D eigenvalues:
    \begin{equation}\lambda_{\mathbf{j}} = \sum_{k=1}^d \lambda_{j_k} = \sum_{k=1}^d \left(\frac{j_k\pi}{L_k+2\delta_k}\right)^2 = \Theta(\norm{\mathbf{j}}_2^2).\end{equation}
    
    \item \textbf{Collocation Grid and Matrices:} The grid is a tensor product of $N$ points in each dimension, resulting in $M = N^d$ points in $R_0$. The $M \times M$ synthesis matrix $\mat{W}_d$ mapping coefficients to function values is the Kronecker product of the 1D matrices: $\mat{W}_d = \mat{W}_N \otimes \dots \otimes \mat{W}_N$. For the Poisson operator $\mathcal{L}_P = -\Delta$, the associated basis functions are $\phi_{\mathbf{j}} = -\Delta w_{\mathbf{j}} = \lambda_{\mathbf{j}} w_{\mathbf{j}}$. The collocation matrix is thus $\mat{A}_P^{(d)} = \mat{W}_d \mathbf{\Lambda}_d$, where $\mathbf{\Lambda}_d$ is the diagonal matrix of eigenvalues $\lambda_{\mathbf{j}}$.
\end{itemize}

%--------------------------------------------------------------------------------
\subsection{Compounding Instability of the Basis}
%--------------------------------------------------------------------------------

We first establish that the baseline instability identified in Section 2 is amplified in higher dimensions.

\begin{theorem}[Exponential Growth of Condition Number in d-Dimensions]\label{thm:d_dim_instability}
For the $d$-dimensional Poisson operator on $R_0 \subset R$ with extension parameters $\delta_1, \dots, \delta_d$, the condition number of the spectral collocation matrix $\mat{A}_P^{(d)}$ grows exponentially with respect to the number of modes. Specifically, let $\alpha = \alpha(\delta_k, L_k)$ denote the 1D decay rate from Theorem~\ref{thm:poisson_illcond} (assuming uniform extensions $\delta_k = \delta$, $L_k = L$ for simplicity). Then
\begin{equation} \kappa(\mat{A}_P^{(d)}) \ge C e^{\alpha(\delta, L)\, d N}, \end{equation}
where $C > 0$ is independent of $N$. The compounding factor $d$ in the exponent reflects the tensor-product structure of the frame.
\end{theorem}
\begin{proof}
The condition number is bounded below by the spectral norm of the inverse matrix. Using the factorization $\mat{A}_P^{(d)} = \mat{W}_d \mathbf{\Lambda}_d$, we have:
\begin{equation} \kappa(\mat{A}_P^{(d)}) \ge \norm{(\mat{A}_P^{(d)})^{-1}}_2 = \norm{\mathbf{\Lambda}_d^{-1} \mat{W}_d^{-1}}_2 \ge \sigma_{\min}(\mathbf{\Lambda}_d^{-1}) \norm{\mat{W}_d^{-1}}_2. \end{equation}
We analyze the two components:
\begin{enumerate}
    \item The eigenvalues scale polynomially. The smallest singular value of $\mathbf{\Lambda}_d^{-1}$ corresponds to the inverse of the largest eigenvalue, which scales as $\Theta(N^{-2})$.
    \item The frame matrix scales exponentially. The matrix $\mat{W}_d$ is the $d$-fold Kronecker product of the 1D matrix $\mat{W}_N$. The singular values of a Kronecker product are the products of the singular values of the constituent matrices. Thus:
    \begin{equation} \sigma_{\min}(\mat{W}_d) = (\sigma_{\min}(\mat{W}_N))^d. \end{equation}
    From the 1D analysis (Theorem~\ref{thm:poisson_illcond}), we know $\sigma_{\min}(\mat{W}_N) \le C_1 e^{-\alpha(\delta,L)\, N}$. Therefore, the smallest singular value of the $d$-dimensional frame is:
    \begin{equation} \sigma_{\min}(\mat{W}_d) \le (C_1)^d e^{-\alpha(\delta,L)\, d N}. \end{equation}
    Consequently,
    \begin{equation} \norm{\mat{W}_d^{-1}}_2 = 1/\sigma_{\min}(\mat{W}_d) \ge C_2 e^{\alpha(\delta, L)\, d N}. \end{equation}
\end{enumerate}
Combining these, the polynomial decay of the eigenvalues is dominated by the exponential growth of the inverse frame matrix. Thus, the condition number grows as $e^{\alpha(\delta, L)\, d N}$.
\end{proof}

%--------------------------------------------------------------------------------
\subsection{Pre-Asymptotic Structural Dichotomy in d-Dimensions}
%--------------------------------------------------------------------------------

Despite the baseline instability of the basis, the pre-asymptotic structural dichotomy remains valid. The inverse of the convection-diffusion operator retains its highly asymmetric structure, which explains the method's practical robustness in multi-dimensional simulations.

\begin{theorem}[Pre-Asymptotic Structural Dichotomy in d-Dimensions]
\label{thm:d-dim-dichotomy}
Let the spectral collocation method be applied to the constant-coefficient convection-diffusion operator $\mathcal{L}_{CD} = -\Delta + \mathbf{k} \cdot \nabla$ in $d$ dimensions with homogeneous Dirichlet boundary conditions and extension parameters $\delta_1, \dots, \delta_d$. Assume the discretization is stabilized so that the interpolation Lebesgue constant $\Lambda_N^{(d)}$ is bounded (or its growth is controlled). Then for any $0 < \alpha < |\mathbf{k}|$, the inverse of the discrete physical-space operator, $(\mat{L}_{CD,N}^{(d)})^{-1}$, is numerically highly asymmetric with entries that decay exponentially in the upstream direction. Specifically, there exists a constant $C > 0$, independent of $N$, such that
\begin{equation}|((\mat{L}_{CD,N}^{(d)})^{-1})_{ij}| \le C\, N^{d/2} \Lambda_N^{(d)}\, \exp\bigl(-\alpha \max(0,\, -\hat{\mathbf{k}} \cdot (\mathbf{x}_i - \mathbf{x}_j))\bigr),\end{equation}
where $\hat{\mathbf{k}} = \mathbf{k}/|\mathbf{k}|$ is the unit convection direction and $\Lambda_N^{(d)}$ is the Lebesgue constant of the $d$-dimensional interpolation operator. In the downstream direction ($\hat{\mathbf{k}} \cdot (\mathbf{x}_j - \mathbf{x}_i) \le 0$), the exponential decay does not apply, and the entries form an $O(1)$ dense block bounded by $C N^{d/2} \Lambda_N^{(d)}$.
\end{theorem}
\begin{proof}
The proof follows the same Combes--Thomas strategy as the 1D case (Theorem~\ref{thm:exp_decay_cd}), with the key modification that the Sobolev embedding $H^1 \hookrightarrow L^\infty$ fails for $d \ge 2$. We address this in three steps.

\noindent\textbf{1. Structure of the Continuous Green's Function.}
For the constant-coefficient operator $\mathcal{L}_{CD} = -\Delta + \mathbf{k} \cdot \nabla$ in $\mathbb{R}^d$, the fundamental solution $\Phi(\mathbf{x})$ can be derived directly via the substitution $\Phi(\mathbf{x}) = \exp(\frac{1}{2}\mathbf{k} \cdot \mathbf{x}) v(\mathbf{x})$. Substituting this into $(-\Delta + \mathbf{k} \cdot \nabla)\Phi = \delta(\mathbf{x})$ transforms the problem into the modified Helmholtz equation $(-\Delta + \frac{1}{4}|\mathbf{k}|^2)v = \delta(\mathbf{x})$. The far-field behavior of $v(\mathbf{x})$ is governed by the modified Bessel function of the second kind, which yields the asymptotic expansion $v(\mathbf{x}) \sim C |\mathbf{x}|^{(1-d)/2} \exp(-\frac{1}{2}|\mathbf{k}||\mathbf{x}|)$. Transforming back provides the far-field behavior of the convection-diffusion fundamental solution: $\Phi(\mathbf{x}) \sim C |\mathbf{x}|^{(1-d)/2} \exp\bigl(\tfrac{1}{2}\mathbf{k} \cdot \mathbf{x} - \tfrac{1}{2}|\mathbf{k}||\mathbf{x}|\bigr)$. Since $\mathbf{k} \cdot \mathbf{x} \le |\mathbf{k}||\mathbf{x}|$, the exponent is strictly non-positive. This produces strong exponential decay upstream ($\hat{\mathbf{k}} \cdot \mathbf{x} < 0$), while avoiding exponential growth downstream ($\hat{\mathbf{k}} \cdot \mathbf{x} > 0$), where it instead decays algebraically. On the bounded domain $R_0$, the continuous Green's function $G_{CD}^{(d)}(\mathbf{x}, \mathbf{y})$ possesses an integrable singularity at $\mathbf{x}=\mathbf{y}$, but away from the singularity it decays strictly in the upstream direction bounded by $C \exp(-\alpha \max(0, -\hat{\mathbf{k}} \cdot (\mathbf{x} - \mathbf{y})))$ for any $\alpha < |\mathbf{k}|$. Because the discrete space $V_{N,d}$ is comprised of band-limited analytic functions, the discrete Green's function cannot support true singularities. Its elements are instead regularized and remain uniformly bounded via the maximum-norm stability proven below in Step 2.

\noindent\textbf{2. Algebraic Bound on the Inverse.}
Unlike the 1D case, we cannot rely on the continuous Sobolev embedding $H^1(R_0) \hookrightarrow L^\infty(R_0)$, which fails for $d \ge 2$. Instead, we establish a direct algebraic bound on the inverse operator in Lemma~\ref{lem:algebraic_bound_d_dim} below, by applying a polynomial inverse inequality. The resulting bound takes the form $\norm{(\mat{L}_{CD,N}^{(d)})^{-1}}_\infty \le C_{inv} \cdot N^{d/2} \Lambda_N^{(d)}$.

\noindent\textbf{3. Combes--Thomas Argument.}
With the algebraic bound established, we apply the same weighted stability argument as in Lemma~\ref{lem:combes_thomas} (the 1D case). Fix a column index $j$ and define the directional weight matrix $\mat{D}_{-\alpha}$ with $(\mat{D}_{-\alpha})_{ii} = e^{-\alpha \hat{\mathbf{k}} \cdot (\mathbf{x}_i - \mathbf{x}_j)}$. The conjugated operator $\widetilde{\mat{L}} = \mat{D}_{-\alpha} \mat{L}_{CD,N}^{(d)} \mat{D}_{-\alpha}^{-1}$ corresponds at the continuous level to the modified operator $\widetilde{\mathcal{L}} = -\Delta + (\mathbf{k} - 2\alpha\hat{\mathbf{k}}) \cdot \nabla + \alpha(|\mathbf{k}| - \alpha)$, which has a strictly positive zeroth-order term for $\alpha < |\mathbf{k}|$. The continuous conjugated operator is strictly coercive, so the discrete conjugated scheme inherits unconditional $H^1$ stability for sufficiently large $N$. Applying Lemma~\ref{lem:algebraic_bound_d_dim} to the conjugated system gives $\norm{\widetilde{\mat{L}}^{-1}}_\infty \le \tilde{C}_{inv} \cdot N^{d/2} \Lambda_N^{(d)}$, from which the entry-wise bound follows:
\begin{equation} |((\mat{L}_{CD,N}^{(d)})^{-1})_{ij}| = |(\mat{D}_{-\alpha}^{-1} \widetilde{\mat{L}}^{-1} \mat{D}_{-\alpha})_{ij}| \le e^{\alpha \hat{\mathbf{k}} \cdot (\mathbf{x}_i - \mathbf{x}_j)} \norm{\widetilde{\mat{L}}^{-1}}_\infty \le \tilde{C}_{inv}\, N^{d/2} \Lambda_N^{(d)}\, e^{\alpha \hat{\mathbf{k}} \cdot (\mathbf{x}_i - \mathbf{x}_j)}. \end{equation}
Since $\hat{\mathbf{k}} \cdot (\mathbf{x}_i - \mathbf{x}_j)$ can be positive (downstream), the conjugated bound is only useful for the upstream direction. For the downstream direction, we rely on the unweighted stability bound $\norm{(\mat{L}_{CD,N}^{(d)})^{-1}}_\infty \le C_{inv}\, N^{d/2} \Lambda_N^{(d)}$, which mathematically confirms that the discrete Green's function plateaus to an $O(1)$ constant block, correctly matching the physics of multi-dimensional convection where downstream information propagates without attenuation.
\end{proof}

\begin{lemma}[Algebraic Bound on the Inverse Operator in $d$-Dimensions]
\label{lem:algebraic_bound_d_dim}
For a stabilized spectral collocation method applied to a strictly coercive convection-diffusion operator $\mathcal{L}_{CD}$ in $d$ dimensions, the inverse physical-space operator satisfies:
\begin{equation} \norm{(\mat{L}_{CD,N}^{(d)})^{-1}}_\infty \le C N^{d/2} \Lambda_N^{(d)}, \end{equation}
where $C$ is a constant independent of $N$, and $\Lambda_N^{(d)}$ is the Lebesgue constant of the stabilized interpolation operator.
\end{lemma}

\begin{proof}
Let $\vect{u} = (\mat{L}_{CD,N}^{(d)})^{-1} \vect{f}$, which corresponds to the discrete solution $u_N \in V_{N,d}$ for an interpolated forcing function $f_I$. By the $H^1$-stability of the stabilized scheme (as established in Lemma~\ref{lem:stability_proof}), the solution energy is bounded by the data:
\begin{equation}
\norm{u_N}_{H^1(R_0)} \le C_{num} \norm{f_I}_{L^2(R_0)}.
\end{equation}

In higher dimensions ($d \ge 2$), we cannot use the Sobolev embedding $H^1 \hookrightarrow L^\infty$ to bound the pointwise values. Instead, we directly apply the polynomial inverse inequality. Because the extended-domain frame is stabilized, the $L^2$ norm of any polynomial $u_N \in V_{N,d}$ on the extended domain $R$ is bounded by its norm on the physical sub-domain $R_0$: $\norm{u_N}_{L^2(R)} \le C_{frame} \norm{u_N}_{L^2(R_0)}$. The standard multi-dimensional trigonometric inverse inequality then yields:
\begin{equation} \norm{u_N}_{L^\infty(R_0)} \le \norm{u_N}_{L^\infty(R)} \le C_{inv} N^{d/2} \norm{u_N}_{L^2(R)} \le C_{inv} C_{frame} N^{d/2} \norm{u_N}_{L^2(R_0)}. \end{equation}

Combining these bounds, and noting that the stabilized Lebesgue constant limits the data interpolation error $\norm{f_I}_{L^2(R_0)} \le C' \Lambda_N^{(d)} \norm{\vect{f}}_\infty$, we obtain:
\begin{equation} \norm{u_N}_{L^\infty(R_0)} \le C N^{d/2} \Lambda_N^{(d)} \norm{\vect{f}}_\infty. \end{equation}
Taking the supremum over all forcing vectors gives the stated bound on the inverse matrix norm.
\end{proof}

\begin{remark}[Impact of the $H^1$ Embedding Failure in Higher Dimensions]
In 1D, the proof seamlessly transitioned from $H^1$ stability to uniform $O(1)$ $\ell_\infty$ boundedness via the Sobolev embedding theorem ($H^1(I) \hookrightarrow L^\infty(I)$). This critical embedding fails for dimensions $d \ge 2$. By relying on the polynomial inverse inequality instead, we necessarily incur an algebraic $N^{d/2}$ growth factor in the matrix norm. However, because $N^{d/2}$ grows only algebraically while the Combes--Thomas spatial attenuation is exponential ($e^{-\alpha x}$), the highly asymmetric structural dichotomy of the Green's function is strictly preserved, and the exponential decay easily overpowers the algebraic penalty at moderate distances. Furthermore, because this bound relies only on standard $H^1$ coercivity, it holds unconditionally for any well-posed convection-diffusion operator, without artificial restrictions on the divergence of expansive convection fields.
\end{remark}

%================================================================================
\section{Numerical Validation}
%================================================================================

We now present numerical experiments to validate the theoretical framework. We analyze the conditioning of the square system ($M=N$) to verify the baseline instability, and we visualize the discrete Green's function using a stabilized overdetermined solver to confirm the structural dichotomy.

%--------------------------------------------------------------------------------
\subsection{Baseline Instability}
%--------------------------------------------------------------------------------
To empirically verify the asymptotic analysis presented in Section 2 and the pre-asymptotic stabilization predicted in Section 3, we analyze the stability of the square system ($M=N$). The numerical experiments were conducted using the basis recombination method to enforce boundary conditions. It is important to note that while recombination algebraically enforces boundary values, it operates within the same underlying Fourier extension frame. Therefore, the asymptotic conditioning properties analyzed in Sections 2 and 3 remain governing factors.

While Theorem~\ref{thm:poisson_illcond} characterizes the instability via the matrix condition number $\kappa(\mat{A}_P)$, we verify this experimentally by computing the Lebesgue constant $\Lambda_N$. As established in Remark~\ref{rem:lebesgue_condition_link}, $\Lambda_N$ acts as a lower bound for the algebraic instability; if $\Lambda_N$ grows exponentially, the condition number of the linear system must also grow exponentially. This provides a robust, solver-independent metric for the underlying frame instability.

For the Poisson problem, we set the physical domain length to $L=\pi$ and the exact solution to $u(x) = \sin(\pi x/L)e^x$. For the convection-diffusion problem, we set $L=1.0$ with a strong convection coefficient $k=10$, using the exact solution $u(x) = \sin(2\pi x/L)$.

The results, summarized in Figure \ref{fig:validation_1d}, confirm the theoretical dichotomy:

\begin{itemize}
    \item \textbf{Poisson Operator (Dashed Lines):} As predicted by Theorem~\ref{thm:poisson_illcond} (via the link in Remark~\ref{rem:lebesgue_condition_link}), the Lebesgue constant $\Lambda_N$ grows exponentially with $N$ regardless of the extension parameter $\delta$. While increasing $\delta$ reduces the slope of the growth (improving the constant), it does not alter the fundamental exponential instability. Consequently, the $L^2$ error (Figure \ref{fig:validation_1d}b) hits an accuracy floor (around $10^{-5}$ to $10^{-8}$ for $\delta \leq 1$) and diverges for large $N$. In contrast, the Convection-Diffusion operator delays this saturation significantly. For larger extension parameters (e.g., $\delta=4.0$), it maintains spectral convergence down to machine precision ($\approx 10^{-14}$), whereas the Poisson operator stagnates near $10^{-10}$ even for large $\delta$.
    
    \item \textbf{Convection-Diffusion Operator (Solid Lines):} In contrast, the convection-diffusion operator with $k=10$ maintains a lower Lebesgue constant throughout the pre-asymptotic regime. This stability allows the method to achieve near-machine precision accuracy (errors $\approx 10^{-14}$) before the asymptotic frame instability eventually dominates. This validates Remark~\ref{rem:pre_asymptotic}, demonstrating that the highly asymmetric inverse structure effectively suppresses the frame's ill-conditioning for practical discretizations.
\end{itemize}

\begin{figure}[ht!]
    \centering
    \includegraphics[width=\textwidth]{validation_1d_summary.png}
    \caption{Baseline instability. (a) The Lebesgue constant $\Lambda_N$ grows exponentially for both operators, confirming the underlying frame instability. (b) However, the Convection-Diffusion system maintains a Lebesgue constant orders of magnitude lower than Poisson in the pre-asymptotic regime.}
    \label{fig:validation_1d}
\end{figure}



%--------------------------------------------------------------------------------
\subsection{Structural Dichotomy: The Green's Function}
%--------------------------------------------------------------------------------

To visualize the intrinsic structure of the inverse operators (Theorem~\ref{thm:exp_decay_cd}) independently of the frame instability, we compute the numerical Green's function. We utilize an overdetermined Least-Squares spectral method with explicit boundary constraints to ensure the solver itself remains robust.

\textbf{Problem Setting:} We solve for the response to a Gaussian source centered at $x_c = L/2$ on the physical domain $I = [0, \pi]$:
\begin{equation}
\mathcal{L}[u](x) = \exp\left(-\frac{(x-x_c)^2}{2\sigma^2}\right), \quad x_c = \pi/2, \quad \sigma = 0.1.
\end{equation}
The computational domain is extended by $\delta=1.0$. To isolate the structural properties, we employ a stable frame with $N=60$ modes and enforce the boundary conditions $u(0)=u(\pi)=0$ via high-weight constraints in the least-squares system. We compare:
\begin{enumerate}
    \item \textbf{Poisson:} $\mathcal{L}_P = -d^2/dx^2$.
    \item \textbf{Convection-Diffusion:} $\mathcal{L}_{CD} = -d^2/dx^2 + k(d/dx)$, with a moderate convection coefficient $k=5.0$.
\end{enumerate}

Figure \ref{fig:greens_function} plots the computed response $u(x)$.
\begin{itemize}
    \item The \textbf{Poisson} response (Blue) is a global, symmetric arch. In the log-scale plot (Panel b), the decay away from the peak is slow and algebraic, indicating a numerically dense inverse matrix.
    \item The \textbf{Convection-Diffusion} response (Orange) is an asymmetric, localized spike. Panel (b) reveals a distinct linear slope on the logarithmic scale upstream of the source, confirming the \textbf{exponential decay} predicted by our theory. Downstream of the source, the response plateaus to an $O(1)$ constant before dropping at the boundary, consistent with the theoretical asymmetry. This structural asymmetry effectively filters the global coupling of errors inherent to the ill-conditioned basis.
\end{itemize}

\begin{figure}[ht!]
    \centering
    \includegraphics[width=\textwidth]{Figure3_Final_ConstraintMethod.png}
    \caption{Visual verification of the Structural Dichotomy (Theorem~\ref{thm:exp_decay_cd}).
    (a) \textbf{Linear Scale:} The Poisson response (Blue) is a global symmetric arch, indicating that the inverse matrix is dense; information at the center affects the entire domain. The Convection-Diffusion response (Orange) is a localized, asymmetric pulse.
    (b) \textbf{Log Scale:} The upstream side of the Convection-Diffusion response exhibits a straight linear slope, confirming the \textbf{exponential decay} predicted by our theory, while the downstream side plateaus. This structural asymmetry isolates numerical errors, explaining the method's robustness in convection-dominated regimes.}
    \label{fig:greens_function}
\end{figure}


%--------------------------------------------------------------------------------
\subsection{2D Validation}
%--------------------------------------------------------------------------------
We extend the validation to two dimensions to confirm that the structural robustness is preserved in tensor-product domains, as discussed in Section 4. The solver utilizes a tensor-product basis formed by the 1D eigenfunctions, $w_{\mathbf{j}}(\mathbf{x}) = w_{j_1}(x_1)w_{j_2}(x_2)$, on the domain $\Omega = [0,1]^2$.

We compare three test cases:
\begin{enumerate}
    \item The Poisson equation $\mathcal{L}_P = -\Delta$.
    \item Constant-coefficient convection-diffusion with $\mathbf{k} = (10, 10)^T$.
    \item Variable-coefficient convection-diffusion with $\mathbf{k}(x,y) = (5\cos(\pi x), 5\cos(\pi y))^T$.
\end{enumerate}

Figure  \ref{fig:validation_2d} presents the $L^2$ convergence history. The 2D Poisson operator (Figure  \ref{fig:validation_2d}a) inherits the instability of the 1D case; the solver fails to reach high accuracy for small $\delta$ and diverges rapidly as $N$ increases, confirming the compounding instability described in Theorem \ref{thm:d_dim_instability}.

Conversely, both the constant (Figure \ref{fig:validation_2d}b) and variable (Figure \ref{fig:validation_2d}c) coefficient convection-diffusion cases exhibit robust spectral convergence. The stabilizing effect of the convection term is even more pronounced in 2D, consistent with the robust structural dichotomy of the convection-diffusion operator. The variable coefficient case demonstrates that this stabilization is robust unconditionally for any well-posed flow even when the convection field varies spatially, provided the local P\'eclet number remains sufficiently high in the pre-asymptotic regime.

\begin{figure}[ht!]
    \centering
    \includegraphics[width=\textwidth]{validation_2d_summary.png}
    \caption{2D Convergence. (a) Poisson. (b) Constant Coeff CD. (c) Variable Coeff CD. The method converges spectrally, demonstrating that the structural robustness extends to tensor-product domains.}
    \label{fig:validation_2d}
\end{figure}

%================================================================================
\section{Conclusion}
%================================================================================

In this work, we have presented a stability analysis of the extended-domain spectral collocation method, reconciling its theoretical asymptotic instability with its robust pre-asymptotic behavior in practice. By shifting the focus from the redundancy of the Fourier extension frame to the structural properties of the discrete inverse operators, we have established a new mathematical framework for understanding this method.

Our analysis confirms the baseline asymptotic failure: the spectral collocation method is asymptotically unstable for both the Poisson and convection-diffusion operators, driven by an exponentially growing Lebesgue constant intrinsic to the trial space. However, the central contribution of this paper is the proof that the method's effectiveness in convection-dominated problems is a distinct pre-asymptotic phenomenon, governed by the physics of the differential operator.

We have proven that the non-self-adjoint nature of the convection-diffusion operator induces a numerically highly asymmetric discrete Green's function. Its entries decay exponentially in the upstream direction, while downstream they plateau to a stable, $O(1)$ constant, consistent with the underlying physics. In contrast, the self-adjoint Poisson operator generates a numerically dense inverse with global algebraic decay. We have shown that the robustness of the unstabilized square system ($M=N$) is rooted in this structural asymmetry. While both operators possess an $O(\Lambda_N)$ inverse norm bound, the dense Poisson inverse amplifies the frame's instability globally, whereas the asymmetric convection-diffusion inverse localizes and suppresses it.

Furthermore, we extended this structural dichotomy to $d$-dimensional tensor-product domains. While the frame instability compounds exponentially with dimension, we demonstrated that the discrete Green's function retains its defining asymmetry: exponential decay strictly upstream and a stable downstream plateau. In higher dimensions, the failure of the Sobolev embedding introduces an algebraic $N^{d/2}$ growth penalty, but this penalty is strictly overpowered by the exponential spatial attenuation of the continuous physics. 

Together, these findings provide a solid theoretical foundation for the extended-domain method. They offer practitioners a clear understanding of why the method succeeds in pre-asymptotic regimes, how stabilization interacts with the underlying physics, and what the ultimate theoretical limitations are. Future work could extend this structural analysis of the inverse operator to problems in more complex, non-rectangular geometries.

\section*{Acknowledgement}
I acknowledge the financial support from the National Science and Technology Council of Taiwan (NSTC 111-2221-E-002-053-MY3).

%    Bibliographies can be prepared with BibTeX using amsplain,
%    amsalpha, or (for "historical" overviews) natbib style.
%    Insert the bibliography data here.



\bibliography{ref}

\end{document}


